\begin{definition} \label{definition:affine_variety_over_a_field}
Let $k$ be a \CrefAndHyperrefIfExist{definition:field}{field}. An \hldef{affine variety over $k$} is an \CrefAndHyperrefIfExist{definition:closed_subscheme_of_a_scheme}{closed subscheme} of the affine space $\mathbb{A}^n_k$ (\Cref{definition:affine_space_of_dimension_n_over_a_scheme}) for some integer $n \geq 0$.


Equivalently, an affine variety over $k$ is a scheme $X$ that is isomorphic to the spectrum of a finitely generated $k$-algebra:
$$\hlin{X \cong \operatorname{Spec}(A),}$$
where $A = k[T_1, \dots, T_n]/I$ for some ideal $I \subseteq k[T_1, \dots, T_n]$.

As is the case with the general notion of \CrefAndHyperrefIfExist{definition:variety_over_a_field}{variety}, modern definitions of affine varieties often also include the assumption that the variety be \CrefAndHyperrefIfExist{definition:reduced_scheme}{reduced}. Unless otherwise specified, we assume that affine varieties are reduced.


An affine variety over $k$ is a special case of a \CrefAndHyperrefIfExist{definition:quasi_projective_variety_over_a_field}{quasi-projective variety over $k$}, since any closed subscheme of $\mathbb{A}^n_k$ can be viewed as a locally closed subscheme of $\mathbb{P}^n_k$.
\end{definition}
